\begin{pastexam}[label={px:ncu-107-dsa-07}]{中央資工所 107 學年度 第 7 題}
（單選；每題 4 分，答錯倒扣一分）

Consider a variation of the towers of Hanoi problem. We no longer assume that all the disks are
initially on one peg. They may be arbitrarily distributed among the three pegs, say pegs A, B, and
C, as long as they are ordered in decreasing sizes on each peg. The purpose of the puzzle remains
to move all disks to one specified peg, under the same constraints as the original problem, with
as few moves as possible. We use array d[\,] to store such a distribution. For example, if
d[1] = A, d[2] = C, d[3] = C, it means disc 1 (the disc of smallest size) is on peg A, and discs 2
and 3 are on peg C. If peg C is the destination peg, then only 1 move is needed. If peg A is the
destination peg, then it needs 6 moves. Now, given d[1] = A, d[2] = C, d[3] = C, d[4] = B,
d[5] = B, d[6] = A, how many moves are needed to move all 6 disks to peg C.

(a)~35 \qquad (b)~36 \qquad (c)~37 \qquad (d)~38 \qquad (e)~39.
\end{pastexam}

\begin{pastsoln}
答案是 (e) 39。

令 $f(k, T)$ 是把盤子 $1..k$（最小的 $k$ 個）全部集中到柱子 $T$ 最少要幾步，其他較大的盤子不動。
看最大的盤子 $k$：
\begin{itemize}
\item \textbf{盤子 $k$ 已經在 $T$ 上}：它比 $1..k-1$ 都大，放在最底下不妨礙任何較小的盤子，不必動它。
$f(k, T) = f(k-1, T)$。
\item \textbf{盤子 $k$ 在別的柱子 $P \ne T$ 上}：它一定要搬到 $T$。 搬它的那一步，$1..k-1$ 不能壓在它
上面，也不能在 $T$ 上，只能全部疊在第三根柱子 $X$（$X \ne P, T$）。 所以先花 $f(k-1, X)$ 步把它們集中到
$X$，搬盤子 $k$ 一步，再把 $X$ 上疊好的 $k - 1$ 個盤子整疊搬到 $T$，這是標準的河內塔，要
$2^{k-1} - 1$ 步（資結書第~6.4~節）：
\[
f(k, T) = f(k-1, X) + 1 + (2^{k-1} - 1) = f(k-1, X) + 2^{k-1}.
\]
\end{itemize}
$f(0, \cdot) = 0$。 這就是本題的遞迴式（資結書第~2.1~節的寫法：由演算法的結構寫出步數）。 盤子 $k$ 多搬幾次
不會更省：每搬一次，較小的盤子都得先整疊讓開，只會多花步數；後面整疊搬 $k - 1$ 個盤子也不能少於
$2^{k-1} - 1$ 步。

\textbf{先驗證題目的例子}。 $d = (A, C, C)$：
\begin{itemize}
\item 目標 C：盤子 $3$、$2$ 都在 C 上，$f(3, C) = f(2, C) = f(1, C)$；盤子 $1$ 在 A，$f(1, C) = f(0, \cdot) + 2^0 = 1$。 和題目說的 $1$ 步相同。
\item 目標 A：盤子 $3$ 在 C，$X = B$，$f(3, A) = f(2, B) + 4$。 盤子 $2$ 在 C，$X = A$，
$f(2, B) = f(1, A) + 2$；盤子 $1$ 已經在 A，$f(1, A) = 0$。 所以 $f(3, A) = 0 + 2 + 4 = 6$，也和題目相同。
\end{itemize}

\textbf{本題}。 $d = (A, C, C, B, B, A)$，目標 C，由大盤子往小盤子算：
\begin{center}
\begin{tabular}{clll}
\toprule
要算的 & 盤子 $k$ 在 & 情形 & 結果 \\
\midrule
$f(6, C)$ & A & 不在 C，$X = B$ & $f(5, B) + 2^5$ \\
$f(5, B)$ & B & 已在 B & $f(4, B)$ \\
$f(4, B)$ & B & 已在 B & $f(3, B)$ \\
$f(3, B)$ & C & 不在 B，$X = A$ & $f(2, A) + 2^2$ \\
$f(2, A)$ & C & 不在 A，$X = B$ & $f(1, B) + 2^1$ \\
$f(1, B)$ & A & 不在 B & $f(0, \cdot) + 2^0 = 1$ \\
\bottomrule
\end{tabular}
\end{center}
往回代：$f(2, A) = 1 + 2 = 3$，$f(3, B) = 3 + 4 = 7$，$f(4, B) = f(5, B) = 7$，
$f(6, C) = 7 + 32 = 39$。

實際的搬法：盤子 $1$ A$\to$B，盤子 $2$ C$\to$A，盤子 $1$ B$\to$A（$3$ 步，$1, 2$ 疊在 A）；盤子 $3$ C$\to$B，
再把 A 上的 $1, 2$ 搬到 B（$1 + 3 = 4$ 步，$1..5$ 都在 B）；盤子 $6$ A$\to$C（$1$ 步）；最後把 B 上的 $5$ 個
盤子搬到 C（$2^5 - 1 = 31$ 步）。 合計 $3 + 4 + 1 + 31 = 39$。

其他選項：(a)--(d) 都比 $39$ 少。 盤子 $6$ 要從 A 搬到 C，搬之前 $1..5$ 必須全在 B，之後 $5$ 個盤子從 B 搬到
C 至少 $31$ 步；把 $1..5$ 集中到 B 至少 $f(5, B) = 7$ 步，所以至少 $7 + 1 + 31 = 39$ 步。
\end{pastsoln}
